\chapter{AI Integration}
\label{ch:ai}

\haira{} provides first-class AI integration for code generation. AI blocks allow developers to describe intent in natural language, with implementations generated at compile time.

\section{Philosophy}

\begin{itemize}
    \item \textbf{Explicit invocation} -- AI is only used when explicitly requested via \keyword{ai} blocks
    \item \textbf{Compile-time generation} -- All AI processing happens during compilation
    \item \textbf{Local-first} -- Prefer local models (llama.cpp, Ollama) over cloud APIs
    \item \textbf{Deterministic builds} -- Generated code is cached for reproducibility
    \item \textbf{Type-safe output} -- Generated code must match declared signatures
\end{itemize}

\section{AI Block Syntax}

\begin{lstlisting}
ai function_name(param1: Type1, param2: Type2) -> ReturnType {
    Natural language description of what the function should do.
    Can span multiple lines.
    Include examples and constraints as needed.
}
\end{lstlisting}

\subsection{Basic Example}

\begin{lstlisting}
import "io"

struct User {
    name: string
    email: string
    created_at: int
}

ai format_user(user: User) -> string {
    Format the user information into a human-readable string.
    Include the name, email, and how long ago they registered.
    Example: "Alice (alice@example.com) - joined 3 days ago"
}

fn main() {
    user = User{
        name: "Alice",
        email: "alice@example.com",
        created_at: 1709251200
    }
    io.println(format_user(user))
}
\end{lstlisting}

\subsection{Complex Example}

\begin{lstlisting}
struct Transaction {
    id: string
    amount: float
    currency: string
    timestamp: int
    category: string
}

ai analyze_spending(transactions: []Transaction) -> SpendingReport {
    Analyze the list of transactions and generate a spending report.

    The report should include:
    - Total spending per category
    - Average transaction amount
    - Highest and lowest transactions
    - Month-over-month comparison if data spans multiple months

    Handle edge cases:
    - Empty transaction list should return a report with zero values
    - Mixed currencies should be converted to USD using approximate rates
}

struct SpendingReport {
    by_category: [string:float]
    average_amount: float
    highest: Transaction?
    lowest: Transaction?
    mom_change: float?
}
\end{lstlisting}

\section{AI Backend Configuration}

Configure AI backends in \code{haira.toml}:

\begin{lstlisting}[language={}]
[ai]
# Primary backend (choose one)
backend = "llama.cpp"    # Local llama.cpp
# backend = "ollama"     # Local Ollama server
# backend = "anthropic"  # Anthropic API
# backend = "openai"     # OpenAI API

# llama.cpp settings
[ai.llama]
model_path = "~/.haira/models/codellama-7b.gguf"
context_size = 4096
threads = 8
gpu_layers = 32

# Ollama settings
[ai.ollama]
host = "localhost:11434"
model = "deepseek-coder-v2:16b"
timeout = 60

# Anthropic settings
[ai.anthropic]
api_key = "${ANTHROPIC_API_KEY}"
model = "claude-sonnet-4-20250514"

# OpenAI settings
[ai.openai]
api_key = "${OPENAI_API_KEY}"
model = "gpt-4"

# Caching
[ai.cache]
enabled = true
directory = ".haira/ai-cache"
\end{lstlisting}

\subsection{Backend Priority}

When multiple backends are available:

\begin{enumerate}
    \item Try configured primary backend
    \item Fall back to other local backends
    \item Fall back to cloud backends (if configured)
    \item Error if no backend available
\end{enumerate}

\subsection{CLI Overrides}

\begin{lstlisting}[language=bash,numbers=none]
# Use specific backend
haira build main.haira --ai-backend=ollama

# Offline mode (cache only)
haira build main.haira --offline

# Regenerate all AI blocks
haira build main.haira --refresh-ai

# Disable AI (error on ai blocks)
haira build main.haira --no-ai
\end{lstlisting}

\section{Compile-Time Processing}

\subsection{Generation Flow}

\begin{enumerate}
    \item Parser encounters \keyword{ai} block
    \item Extract: function signature, intent description, surrounding context
    \item Generate cache key: \code{sha256(intent + signature + context)}
    \item Check cache for existing generation
    \item If not cached: invoke AI backend with structured prompt
    \item Parse generated code (CIR format)
    \item Validate against declared signature
    \item Insert into AST as regular function
    \item Cache result for future builds
\end{enumerate}

\subsection{Context Awareness}

AI blocks have access to:

\begin{itemize}
    \item Function signature (parameters, return type)
    \item Type definitions used in signature
    \item Imported modules and their public interfaces
    \item Other functions in the same file
\end{itemize}

\begin{lstlisting}
struct User {
    name: string
    posts: []Post
}

struct Post {
    title: string
    likes: int
}

// AI can see User and Post definitions
ai get_popular_posts(user: User, min_likes: int) -> []Post {
    Return all posts from the user that have at least min_likes likes.
    Sort by likes in descending order.
}
\end{lstlisting}

\section{CIR: Canonical Intermediate Representation}

AI backends output CIR, a JSON-based intermediate format that represents \haira{} code.

\subsection{CIR Structure}

\begin{lstlisting}[language={}]
{
  "version": "1.0",
  "body": [
    {
      "type": "statement",
      "kind": "variable_decl",
      "name": "result",
      "value": { ... }
    },
    {
      "type": "statement",
      "kind": "return",
      "value": { ... }
    }
  ]
}
\end{lstlisting}

See Chapter~\ref{ch:cir} for the complete CIR specification.

\subsection{Why CIR?}

\begin{itemize}
    \item \textbf{Structured output} -- Easier for AI to generate correctly
    \item \textbf{Validation} -- Can validate structure before code generation
    \item \textbf{Language evolution} -- CIR can remain stable as syntax changes
    \item \textbf{Debugging} -- Easier to inspect and debug than raw code
\end{itemize}

\section{Caching and Reproducibility}

\subsection{Cache Key Generation}

\begin{lstlisting}[language={}]
cache_key = sha256(
    intent_description +
    function_signature +
    type_definitions +
    haira_version +
    ai_backend_version
)
\end{lstlisting}

\subsection{Cache Location}

\begin{lstlisting}[language={}]
.haira/
  ai-cache/
    ab12cd34.json     # Cached CIR
    ab12cd34.meta     # Metadata (timestamps, backend used)
\end{lstlisting}

\subsection{Cache Commands}

\begin{lstlisting}[language=bash,numbers=none]
# View cache statistics
haira ai cache-stats

# Clear cache
haira ai cache-clear

# Export cache (for CI/CD)
haira ai cache-export --output cache.tar.gz

# Import cache
haira ai cache-import cache.tar.gz
\end{lstlisting}

\section{Best Practices}

\subsection{When to Use AI Blocks}

\textbf{Good use cases:}
\begin{itemize}
    \item Data transformation and formatting
    \item Text processing and parsing
    \item Business logic with clear rules
    \item Utility functions with well-defined behavior
\end{itemize}

\textbf{Poor use cases:}
\begin{itemize}
    \item Security-critical code
    \item Performance-critical hot paths
    \item Code requiring external state or I/O
    \item Complex algorithms with specific requirements
\end{itemize}

\subsection{Writing Good Intent Descriptions}

\begin{lstlisting}
// BAD: Too vague
ai process(data: []int) -> int {
    Process the data.
}

// GOOD: Specific and detailed
ai calculate_median(numbers: []int) -> float {
    Calculate the median value of the given numbers.

    For odd-length arrays, return the middle element.
    For even-length arrays, return the average of the two middle elements.

    Handle edge cases:
    - Empty array: return 0.0
    - Single element: return that element as float

    Example:
    - [1, 3, 5] -> 3.0
    - [1, 2, 3, 4] -> 2.5
}
\end{lstlisting}

\subsection{Testing AI-Generated Code}

\begin{lstlisting}
ai parse_date(s: string) -> (Date, Error?) {
    Parse a date string in various formats:
    - "2024-03-15"
    - "March 15, 2024"
    - "15/03/2024"
    Return error for invalid formats.
}

#[test]
fn test_parse_date() {
    // Test ISO format
    date, err = parse_date("2024-03-15")
    assert(err == nil)
    assert(date.year == 2024)
    assert(date.month == 3)
    assert(date.day == 15)

    // Test invalid format
    _, err = parse_date("not a date")
    assert(err != nil)
}
\end{lstlisting}

\section{Error Handling}

\subsection{Compilation Errors}

\begin{lstlisting}[language={}]
error[AI001]: AI generation failed
  --> src/main.haira:15:1
   |
15 | ai complex_function(...) -> Result {
   | ^^^ AI backend returned invalid CIR
   |
   = note: Backend error: "Unable to generate implementation"
   = help: Try simplifying the intent description

error[AI002]: Generated code type mismatch
  --> src/main.haira:20:1
   |
20 | ai transform(x: int) -> string {
   | ^^^ Expected return type `string`, got `int`
   |
   = note: AI generated code returning wrong type
   = help: Regenerate with --refresh-ai
\end{lstlisting}

\subsection{Fallback Behavior}

\begin{lstlisting}
// Provide fallback implementation
ai smart_parse(input: string) -> Data {
    Parse the input intelligently, handling various formats.
} fallback {
    // Manual implementation if AI fails
    return basic_parse(input)
}
\end{lstlisting}

\section{AI Block Grammar}

\begin{lstlisting}[language=EBNF]
ai_block = "ai" identifier "(" [ params ] ")" [ "->" type ]
           "{" intent_description "}"
           [ "fallback" block ] ;

intent_description = { any_char } ;  (* Free-form text *)
\end{lstlisting}
